\section{Harmonic quotient metrics and varying backgrounds}\label{sec:quotients}

\subsection{The quotient comparison}
We now prove \eqref{eq:quotient-main}. Let $v$ be a tangent to one
of the smooth pieces, represented at a flat-balanced point $x$.
Section~\ref{subsec:pieces} gives a locally smooth correction
$y=g(x)x$. At the point under consideration choose its value
$g=e^S$ with $\|S\|\le C\delta$.
Differentiating this change of presentation gives $gv$ plus a
complex orbit vector. The latter is removed by the finite minimum.
Similarly, a varying holomorphic gauge in the infinite-dimensional
presentation contributes an exact $D_a$ term, removed by the
harmonic minimum.

More explicitly, Proposition~\ref{prop:positive-form} at $y$ gives
\[
 \inf_{\gamma}\|J_yw-D_a\gamma\|^2
       =\inf_{\xi\in\g_0}\B_{t,y}(w-C_y\xi,w-C_y\xi).
\]
The bounded value of $g$ sends finite orbit directions bijectively
to those at $y$, with norm distortion $1+O(\delta)$.
Apply \eqref{eq:B-comparison} before taking the infimum, and then
set $w=gv$ modulo the vertical term. This proves
\[
 (1-C\delta)\inf_{\xi}\|v-C_x\xi\|^2
 \le\|[v]\|_{L^2,t}^2
 \le(1+C\delta)\inf_{\xi}\|v-C_x\xi\|^2.
\]
No estimate for $Dg$ is required. A large derivative of a local
balancing gauge near a stabilizer change thus does not affect
the tensor constant.

At a flat-balanced point, differentiating $\mu=\zeta$ makes a
balanced tangent perpendicular to Hermitian orbit directions.
Centrality of $\zeta$ makes Hermitian and unitary orbit directions
perpendicular. Hence the full complex-orbit minimum agrees with
the compact K\"ahler quotient metric on such tangents.
Positivity of $\B$ also identifies the infinitesimal kernel on
integrable vectors with $\im C_x$. Together with the actual
coverage and smooth pieces already proved, this completes
Theorem~\ref{thm:main}.

The order of minimization is the one in
\cite[Section~8]{wuebben6}: remove the uniformly coercive
infinite-dimensional directions before minimizing over the
possibly singular finite orbit. Taking an inverse of the finite
orbit Hessian first would lose the stated uniformity.
There is also no assertion about arbitrary untransported tangent
limits. For the weight-one action on $\C$, for example, the quotient
seminorm of $v=1$ is zero at $x\ne0$ and one at $x=0$.

\subsection{Moving harmonic backgrounds}\label{subsec:moving}
The additive parameter error in Theorem~\ref{thm:main} can be
removed by choosing the finite model at the current polarization.

\begin{proposition}\label{prop:moving}
Keep the holomorphic factors $F_i$ fixed and let $\omega_t$ vary
near $\omega_0$. There are smooth determinant-normalized HE metrics
$h_i(t)$ with $h_i(0)=h_{i0}$. Form
\[
 h_{\mathrm{gr},t}=\bigoplus_i h_i(t)\otimes I_{W_i},
\]
and use its harmonic spaces $V_t$, norms, Green operators and full
canonical obstruction $\kappa_t$. In these coordinates all
conclusions of Theorem~\ref{thm:main} hold uniformly with error
$C\|x\|_t^2$. The reference metric is $h_{\mathrm{gr},t}$,
the spectral mass is
\[
 \Vol(\omega_t)\sum_i r_i\tr(\xi_i^*\xi_i),
\]
and the central target is still $\zeta(t)$.
The factors may have unequal slopes for $t\ne0$.
\end{proposition}

\begin{proof}
Each central factor is stable and simple. After separately solving
its scalar determinant equation, its trace-free HE linearization
is invertible. The parameter implicit-function theorem gives the
smooth family $h_i(t)$. HE existence implies polystability, and
simplicity of the fixed holomorphic factor then implies stability.
This argument supplies nearby factor stability without assuming a
uniform chamber-size statement.

The holomorphic Dolbeault complexes do not change, so all cohomology
dimensions are fixed. Smooth elliptic perturbation gives smooth
harmonic projections and Green operators, with uniformly positive
complementary gaps on a smaller parameter neighborhood.
Projection followed by the positive Gram square root gives smooth
orthonormal bases in each factor Ext space. Tensoring with
$\Hom(W_i,W_j)$ makes the trivialization $G$-equivariant.
The zero-form kernel is the same multiplicity algebra: cross-Hom
vanishing is a fixed holomorphic fact, and the scalar endomorphisms
of each stable factor are parallel for every $h_i(t)$.

For this construction use the moving density $dV_t$ and moment
representative
\begin{equation}\label{eq:moving-moment}
 \M_t^{\mathrm{mov}}=i\Lambda_{\omega_t}F^0,
 \qquad Q_{H,t}\text{ defined by the moving }L^2\text{ product}.
\end{equation}
At the split connection the mean curvature is a constant on each
factor block, so
$Q_{H,t}\M_t^{\mathrm{mov}}(0)=0$ exactly.
The fixed-density representative \eqref{eq:curvature-moment}
would multiply this block endomorphism by $dV_t/dV_0$; when slopes
differ that generally has a nonzero complementary projection.
Nevertheless the integrated moment is unchanged, because
\[
 \M_t^{\mathrm{mov}}dV_t
       =iF^0\wedge\frac{\omega_t^{n-1}}{(n-1)!}.
\]
Thus the central target and rank conventions remain those of
\eqref{eq:moment}.

Identify the smoothly moving Hermitian complements with a fixed
Sobolev space. The derivative of the complementary equation at
$t=0$ is invertible, and its small perturbation remains uniformly
invertible. At unequal slopes it need not equal a scalar multiple
of a Dolbeault Laplacian on the whole Hermitian endomorphism bundle;
continuity of the actual derivative suffices.
For each fixed $t$, the linear coefficient derivative of contracted
curvature vanishes on a current harmonic one-form. This follows
from the Chern-connection K\"ahler identities, which express the
contraction in terms of $d_t^*x$ and its adjoint and do not require
the entire mean curvature to be central.
Consequently the complementary solution satisfies
\begin{equation}\label{eq:moving-sigma}
 \sigma(t,0)=0,\qquad D_x\sigma(t,0)=0,\qquad
 \|\sigma(t,x)\|_{C^j}\le C_j\|x\|_t^2 .
\end{equation}
Uniform bounds for second coefficient derivatives on a smaller
product neighborhood give the same constant for all small $t$.

At every $t$, the quadratic Kuranishi term is current-coexact and
the quadratic gauge term is current-exact. Both are orthogonal to
the current harmonic space. The proof of
Proposition~\ref{prop:quadratic-form} therefore gives
$\Omega_{t,x}-\omega_{\mathrm{flat},t}=O(\|x\|_t^2)$.
The complementary minimum has the same quadratic error, by
Proposition~\ref{prop:positive-form} in the current norms.
Parallel equivariance and $d_t^*J_{t,x}=0$ give the same
energy-relative Schur bound with uniformly controlled background
gaps. Lemma~\ref{lem:balancing} in the orthonormal quiver coordinates
then gives $\|S\|\le C\|x\|_t^2$, and the actual metric, spectral and
quotient transfers preserve that error. The product determinant
is already HE, since the sum of its constant factor mean curvatures
is constant.

Finally these moving coordinates retain common local coverage.
Impose moving Coulomb gauge on an integrable operator using the
complement of the unchanged holomorphic zero-form algebra.
Its linearization is the moving positive zero-form Laplacian,
uniformly invertible. The parameter slice theorem gives a common
gauge neighborhood. For an integrable Coulomb representative $a$,
Hodge decomposition gives
\[
 a=H_{1,t}a-d_t^*G_{2,t}(a\wedge a),
\]
so contraction uniqueness identifies it with the moving canonical
representative. Positive-gauge normalization and the reverse
finite correction apply as before. Smooth changes between the
unitary presentations for $h_0$ and $h_{\mathrm{gr},t}$ transfer
nested neighborhoods. All quadratic comparisons remain in the
specified moving harmonic coordinates; no arbitrary nonlinear
coordinate change is used to preserve their order.
\end{proof}

